\section{Direct scalar-critic deployment and independent confirmation}\label{sec:r14-experiment}
This study exercises Theorem~\ref{thm:r14-gradient} without an offline face library or a query-time optimal policy. It complements, rather than pools with, the preceding multistage experiment. A star has two review dates and many possible continuations; its large leaf count is not evidence about increasing tree depth.

\subsection{Primitives, frozen design, and unrefined decisions}
There are 127 equally likely continuations and discount $0.9$, giving root weight one, leaf weight $w=0.9/127$, and total discounted review mass $1.9$. Tiers lie in $[0,1]$ with outside tier $1/2$, $\chi_0=2$, and $\chi_i=2w$. All absolute switching weights equal their review weights. The root tariff is one; leaf tariff coefficients are $w k_i/8$, with fixed integers $k_i$ drawn from $8,\ldots,12$. Balanced rational resource offsets make $1/2$ the reoptimized common-tier comparator in every context. The companion specifies the full generating formula and seed, not only the reduced objective. Profits and certificates below are divided by $1.9$; gradient errors refer to the unnormalized value and the stated unit of $h$.

Training and test contexts have independent coordinates: $h$ and $\alpha$ are uniform on the thousandth grid in $[-1,1]$, and $\lambda$ is uniform on the thousandth grid in $[0.01,0.20]$. Thus every deployment has strictly positive absolute friction. Eight independent training seeds each generate 1,024 contexts and a 192-unit random-feature tanh scalar critic with a fitted ridge readout. Within each seed, value-only and value-gradient fits share the examples and hidden features. Their squared-error objectives differ only by the prespecified weight $0.25$ on the squared $h$-derivative error; the ridge coefficient is $10^{-9}$. The frozen design, coefficients, labels, and raw per-seed observations are retained. This is a prospective computational plan within the revision, not an externally registered clinical or field protocol.

An independent shared cohort of 128 contexts measures unrefined proposals. Table~\ref{tab:r14-quality} reports all 2,048 policies. Mean optimum gain is \RFourteenOptimalGain; value-only and value-gradient proposal regrets are \RFourteenValueRegret\ and \RFourteenGradientRegret. Every proposal passes the exact gain gate; root-capacity projection occurs on \RFourteenCapRate\ of proposals. The derivative-weighted mean is better, but its paired regret difference has a 95\% Student interval $[\RFourteenCILow,\RFourteenCIHigh]$, which includes zero. The unit of replication is the training seed, not the 128 reused contexts. This interval is conditional on that test cohort and does not establish a population-wide derivative-supervision advantage.
\begin{table}[t]
\centering
\caption{Unrefined accepted scalar-critic policies: eight paired training seeds, 128 shared independent contexts, and positive switching friction.}
\label{tab:r14-quality}
\begin{tabular}{lrr}
\toprule
Metric & Value only & Value gradient\\
\midrule
Mean normalized regret & $5.965\times10^{-6}$ & $2.210\times10^{-6}$\\
Mean exact normalized gap & $1.534\times10^{-5}$ & $5.367\times10^{-6}$\\
Mean residual bound & $9.131\times10^{-6}$ & $3.194\times10^{-6}$\\
Mean gradient bound & $2.224\times10^{-5}$ & $7.046\times10^{-6}$\\
Partial-gradient squared error & $1.171\times10^{-5}$ & $3.708\times10^{-6}$\\
Immediately certified at $10^{-7}$ & 17.48\% & 19.04\%\\
Negative gain, static gate & 0.00\% & 0.00\%\\
\bottomrule
\end{tabular}
\par\smallskip
\begin{minipage}{0.94\linewidth}
All proposal losses and bounds precede full-objective refinement. The paired primary regret interval includes zero; lower mean error alone does not establish a derivative-training advantage. The exact gap includes rational repair, whereas the analytical bounds describe the exact feasible price response.
\end{minipage}
\end{table}


The observable residual and partial-gradient bounds are small on the profit scale, unlike a condition requiring an unknown optimal face. Exact certification is stricter: only \RFourteenImmediateValue\ and \RFourteenImmediateGradient\ of the respective proposals immediately meet the normalized $10^{-7}$ tolerance. Table~\ref{tab:r14-timing} in the companion charges prediction, capacity repair, rational auditing, and every executed refinement, with cold scalar Newton, Brent, and sorted breakpoints as strong structure-aware comparators. In this environment cold Newton's median complete time is \RFourteenColdTime\ ms and the value-gradient pipeline's is \RFourteenLearnedTime\ ms. The learned pipeline is slower, despite fewer refinement iterations. Its positive result is accurate, accepted \emph{unrefined} deployment and a usable value-gradient theorem, not a latency claim obtained by omitting audits or choosing only a dense baseline.

\subsection{An executed expected-gain guarantee without classical rescue}
After freezing the two critics from the first prespecified training seed, we draw a separate 2,048-context cohort with a different recorded seed. Neither model is refit or selected using this cohort before evaluation. Both policies use the feasible price response and exact static gate, with no optimization of the full query objective. Because $G_\lambda(y)$ is increasing in $h$ and decreasing in $\lambda$ for every $y\in Y$, and its maximum is convex in $\alpha$, the optimum over the entire context box is bounded by two deterministic corners: $h=1$, $\lambda=0.01$, $\alpha\in\{-1,1\}$. Exact conjugate upper bounds at those corners give a uniform normalized gain range $[0,\RFourteenRange]$ for every gated policy.

A one-sided Hoeffding bound and a union bound over $M=2$ frozen policies therefore give simultaneous expected-gain lower bounds with $n=2048$ and $\delta=0.05$. The selected value-gradient policy has empirical gain \RFourteenValidationGain\ and lower bound \RFourteenValidationFloor\ per discounted review. Its average exact suboptimality certificate is \RFourteenValidationGap; this last number is an empirical average, not a confidence bound on expected regret. Both policies have positive gain lower bounds, and their small gain difference is not established as significant. This confirmation concerns the declared fixed synthetic distribution, not arbitrary shift or new primitive families. It does demonstrate positive accepted neural deployment without the classical rescue needed by the earlier poor actors.

\subsection{Structural cross-checks and preservation of adverse evidence}
The direct threshold agrees with 160 general tension linear programs. Across 480 solver comparisons, safeguarded Newton, Brent, bisection, and the breakpoint scan agree; 128 small-instance SLSQP comparisons have maximum value discrepancy \RFourteenDenseError. A further 800 arbitrary-price checks include 106 root-capacity-binding cases and verify both global bounds and their underlying remainder identity. These are implementation tests, not substitutes for the proofs. Timing the explicit threshold and scalar solvers through 32,767 leaves confirms the intended sparse implementation; the raw measurements and their two-date limitation appear in Table~\ref{tab:r14-scaling}.

A separate standard-library verifier reconstructs \RFourteenCertificates\ rational policy certificates and \RFourteenChecks\ exact algebraic checks, including every continuation cap, root equality, absolute switching identity, profit, and conjugate upper bound. The full nonlinear and boundary-inclusive predecessor experiment, its badly shifted actors, unresolved derivative comparison, and slower complete learned pipelines remain unchanged. The new result identifies a tractable institution and learning statistic under which the original integrated research objective succeeds; it does not erase failure in the broader distributions.
